\begin{thebibliography}{99}
\bibitem{source1} AmirMahdi Ahmadinejad, Jonathan A. Kelner, Jack Murtagh, John Peebles, Aaron Sidford, and Salil P. Vadhan. High-precision estimation of random walks in small space. In Sandy Irani, editor, 61st IEEE Annual Symposium on Foundations of Computer Science, FOCS 2020, Durham, NC, USA, November 16-19, 2020, pages 1295–1306. IEEE, 2020. doi:10.1109/ FOCS46700.2020.00123.
\bibitem{source2} Miklós Ajtai. Approximate counting with uniform constant-depth circuits. In Jin-Yi Cai, editor, Advances In Computational Complexity Theory, Proceedings of a DIMACS Workshop, New Jersey, USA, December 3-7, 1990, volume 13 of DIMACS Series in Discrete Mathematics and Theoretical Computer Science, pages 1–20. DIMACS/AMS, 1990. doi:10.1090/DIMACS/ 013/01.
\bibitem{source3} Miklós Ajtai and Michael Ben-Or. A theorem on probabilistic constant depth computations. In Richard A. DeMillo, editor, Proceedings of the 16th Annual ACM Symposium on Theory of Computing, April 30 - May 2, 1984, Washington, DC, USA, pages 471–474. ACM, 1984. doi:10.1145/800057.808715.
\bibitem{source4} Roy Armoni. On the derandomization of space-bounded computations. In International Workshop on Randomization and Approximation Techniques in Computer Science, pages 47–59. Springer, 1998.
\bibitem{source5} László Babai, Noam Nisant, and Márió Szegedy. Multiparty protocols, pseudorandom gen- erators for logspace, and time-space trade-offs. Journal of Computer and System Sciences, 45(2):204–232, 1992.
\bibitem{source6} Mark Braverman, Gil Cohen, and Sumegha Garg. Pseudorandom pseudo-distributions with near-optimal error for read-once branching programs. SIAM Journal on Computing, 49(5):STOC18–242, 2019.
\bibitem{source7} Eshan Chattopadhyay and Jyun-Jie Liao. Optimal error pseudodistributions for read-once branching programs. In 35th Computational Complexity Conference, 2020.
\bibitem{source8} Lijie Chen, William M. Hoza, Xin Lyu, Avishay Tal, and Hongxun Wu. Weighted pseu- dorandom generators via inverse analysis of random walks and shortcutting. In 64th IEEE Annual Symposium on Foundations of Computer Science, FOCS 2023, Santa Cruz, CA, USA, November 6-9, 2023, pages 1224–1239. IEEE, 2023. doi:10.1109/FOCS57990.2023.00072.
\bibitem{source9} Kuan Cheng and William M. Hoza. Hitting sets give two-sided derandomization of small space. In Shubhangi Saraf, editor, 35th Computational Complexity Conference, CCC 2020, July 28-31, 2020, Saarbrücken, Germany (Virtual Conference), volume 169 of LIPIcs, pages 10:1–10:25. Schloss Dagstuhl – Leibniz-Zentrum für Informatik, 2020. doi:10.4230/LIPICS.CCC.2020.10.
\bibitem{source10} Gil Cohen, Dean Doron, Oren Renard, Ori Sberlo, and Amnon Ta-Shma. Error reduction for weighted prgs against read once branching programs. In Valentine Kabanets, editor, 36th Computational Complexity Conference, CCC 2021, July 20-23, 2021, Toronto, Ontario, Canada (Virtual Conference), volume 200 of LIPIcs, pages 22:1–22:17. Schloss Dagstuhl – Leibniz-Zentrum für Informatik, 2021. doi:10.4230/LIPICS.CCC.2021.22.
\bibitem{source11} Gil Cohen, Dean Doron, Oren Renard, Ori Sberlo, and Amnon Ta-Shma. Error reduction for weighted prgs against read once branching programs. In Proceedings of the 36th Computational Complexity Conference, page 1, 2021.
\bibitem{source12} Gil Cohen, Dean Doron, Ori Sberlo, and Amnon Ta-Shma. Approximating iterated multiplic- ation of stochastic matrices in small space. In Barna Saha and Rocco A. Servedio, editors, Proceedings of the 55th Annual ACM Symposium on Theory of Computing, STOC 2023, Or- lando, FL, USA, June 20-23, 2023, pages 35–45. ACM, 2023. doi:10.1145/3564246.3585181.
\bibitem{source13} Joshua Cook. Size bounds on low depth circuits for promise majority. In Nitin Saxena and Sunil Simon, editors, 40th IARCS Annual Conference on Foundations of Software Technology and Theoretical Computer Science, FSTTCS 2020, December 14-18, 2020, BITS Pilani, K K Birla Goa Campus, Goa, India (Virtual Conference), volume 182 of LIPIcs, pages 19:1–19:14. Schloss Dagstuhl – Leibniz-Zentrum für Informatik, 2020. doi:10.4230/LIPICS.FSTTCS.2020.19.
\bibitem{source14} Dean Doron, Edward Pyne, and Roei Tell. Opening up the distinguisher: A hardness to randomness approach for BPL = L that uses properties of BPL. Electron. Colloquium Comput. Complex., TR23-208, 2023. arXiv:TR23-208.
\bibitem{source15} Dean Doron and Roei Tell. Derandomization with minimal memory footprint. In Amnon Ta-Shma, editor, 38th Computational Complexity Conference, CCC 2023, July 17-20, 2023, Warwick, UK, volume 264 of LIPIcs, pages 11:1–11:15. Schloss Dagstuhl – Leibniz-Zentrum für Informatik, 2023. doi:10.4230/LIPICS.CCC.2023.11.
\bibitem{source16} Pooya Hatami and William Hoza. Theory of unconditional pseudorandom generators. Electron. Colloquium Comput. Complex., TR23-019, 2023. arXiv:TR23-019.
\bibitem{source17} William M. Hoza. Better pseudodistributions and derandomization for space-bounded com- putation. In Mary Wootters and Laura Sanità, editors, Approximation, Randomization, and Combinatorial Optimization. Algorithms and Techniques, APPROX/RANDOM 2021, August 16-18, 2021, University of Washington, Seattle, Washington, USA (Virtual Conference), volume
\bibitem{source207} of LIPIcs, pages 28:1–28:23. Schloss Dagstuhl – Leibniz-Zentrum für Informatik, 2021. doi:10.4230/LIPICS.APPROX/RANDOM.2021.28.
\bibitem{source18} William M. Hoza. Recent progress on derandomizing space-bounded computation. Bull. EATCS, 138, 2022. URL: http://eatcs.org/beatcs/index.php/beatcs/artic$\ell$e/view/728.
\bibitem{source19} William M Hoza and David Zuckerman. Simple optimal hitting sets for small-success rl. SIAM Journal on Computing, 49(4):811–820, 2020.
\bibitem{source20} Russell Impagliazzo, Noam Nisan, and Avi Wigderson. Pseudorandomness for network algorithms. In Proceedings of the twenty-sixth annual ACM symposium on Theory of computing, pages 356–364, 1994.
\bibitem{source21} Adam R. Klivans and Dieter van Melkebeek. Graph nonisomorphism has subexponential size proofs unless the polynomial-time hierarchy collapses. SIAM J. Comput., 31(5):1501–1526, 2002. doi:10.1137/S0097539700389652.
\bibitem{source22} Noam Nisan. Pseudorandom generators for space-bounded computation. Comb., 12(4):449–461, 1992. doi:10.1007/BF01305237.
\bibitem{source23} Noam Nisan. RL $\subseteq$ SC. Comput. Complex., 4:1–11, 1994. doi:10.1007/BF01205052.
\bibitem{source24} Noam Nisan and David Zuckerman. Randomness is linear in space. Journal of Computer and System Sciences, 52(1):43–52, 1996.
\bibitem{source25} Aaron (Louie) Putterman and Edward Pyne. Near-optimal derandomization of medium-width branching programs. In Barna Saha and Rocco A. Servedio, editors, Proceedings of the 55th Annual ACM Symposium on Theory of Computing, STOC 2023, Orlando, FL, USA, June 20-23, 2023, pages 23–34. ACM, 2023. doi:10.1145/3564246.3585108.
\bibitem{source26} Edward Pyne, Ran Raz, and Wei Zhan. Certified hardness vs. randomness for log-space. In 64th IEEE Annual Symposium on Foundations of Computer Science, FOCS 2023, Santa Cruz, CA, USA, November 6-9, 2023, pages 989–1007. IEEE, 2023. doi:10.1109/FOCS57990.2023.00061.
\bibitem{source27} Edward Pyne and Salil Vadhan. Pseudodistributions that beat all pseudorandom generators. In 36th Computational Complexity Conference (CCC 2021). Schloss Dagstuhl – Leibniz-Zentrum für Informatik, 2021.
\bibitem{source28} Edward Pyne and Salil P. Vadhan. Pseudodistributions that beat all pseudorandom generators (extended abstract). In Valentine Kabanets, editor, 36th Computational Complexity Conference, CCC 2021, July 20-23, 2021, Toronto, Ontario, Canada (Virtual Conference), volume 200 of LIPIcs, pages 33:1–33:15. Schloss Dagstuhl – Leibniz-Zentrum für Informatik, 2021. doi: 10.4230/LIPICS.CCC.2021.33.
\bibitem{source29} Michael E. Saks and Shiyu Zhou. BPH SPACE(S) $\subseteq$ DSPACE(S 3/2 ). J. Comput. Syst. Sci., 58(2):376–403, 1999. doi:10.1006/JCSS.1998.1616.
\bibitem{source30} Emanuele Viola. On approximate majority and probabilistic time. In 22nd Annual IEEE Conference on Computational Complexity (CCC 2007), 13-16 June 2007, San Diego, California, USA, pages 155–168. IEEE Computer Society, 2007. doi:10.1109/CCC.2007.16.
\bibitem{source31} Emanuele Viola. Randomness buys depth for approximate counting. In Rafail Ostrovsky, editor, IEEE 52nd Annual Symposium on Foundations of Computer Science, FOCS 2011, Palm Springs, CA, USA, October 22-25, 2011, pages 230–239. IEEE Computer Society, 2011. doi:10.1109/FOCS.2011.19.
\end{thebibliography}
